\documentclass[11pt,a4paper]{article}

\usepackage[T1]{fontenc}
\usepackage[utf8]{inputenc}
\usepackage[margin=2.2cm]{geometry}
\usepackage{amsmath,amssymb}
\usepackage{booktabs}
\usepackage{longtable}
\usepackage{array}
\usepackage{xcolor}
\usepackage{fancyhdr}
\usepackage{enumitem}
\usepackage{tcolorbox}
\usepackage[hidelinks]{hyperref}

\definecolor{stampblue}{HTML}{1A3A5C}
\definecolor{pagedteal}{HTML}{4D9C96}
\definecolor{warnorange}{HTML}{C25E0C}
\definecolor{lightgrey}{HTML}{F4F6F8}
\definecolor{rulegrey}{HTML}{CCD4DD}

\pagestyle{fancy}
\fancyhf{}
\lhead{\small\textcolor{stampblue}{Memory Management Module --- Evaluation Answers}}
\rhead{\small QuadraMind}
\cfoot{\thepage}
\renewcommand{\headrulewidth}{0.4pt}

\newcommand{\op}[1]{\textsc{#1}}
\newcommand{\file}[1]{\texttt{\small #1}}

\newtcolorbox{keybox}[1]{colback=lightgrey,colframe=stampblue,
  boxrule=0.6pt,left=6pt,right=6pt,top=5pt,bottom=5pt,
  title={\small\bfseries #1},fonttitle=\color{white}}

\newtcolorbox{warnbox}[1]{colback=white,colframe=warnorange,
  boxrule=0.8pt,left=6pt,right=6pt,top=5pt,bottom=5pt,
  title={\small\bfseries #1},fonttitle=\color{white}}

\title{\vspace{-1.7cm}
  \textcolor{stampblue}{\Large\bfseries On-Chip Memory Management Module}\\[0.35em]
  \large Final-Evaluation Answer Sheet\\[0.2em]
  \normalsize Research problem, gap, novelty, contribution, artefacts,
  methodology,\\ scheme comparison, results, limitations and validation}
\author{\small Memory Subsystem --- QuadraMind DNN Accelerator Simulator}
\date{\small 7 August 2026}

\begin{document}
\maketitle
\thispagestyle{fancy}

\begin{keybox}{One-paragraph summary --- if you get 60 seconds, say this}
A DNN layer's working set does not fit in on-chip SRAM, so tiles must be
staged in and out of a scratchpad. Conventional accelerators do this with a
\emph{runtime} structure --- a page table or a cache tag array --- which costs
a lookup on every access, fetches at page granularity, and places data at
addresses the compiler cannot predict, so it cannot avoid bank conflicts. My
module implements the alternative: because a DNN schedule is fully known at
compile time, a compiler can pre-compute every scratchpad layout and emit only
the \emph{differences} between consecutive layouts as a five-opcode delta
stream, so no runtime tag or translation structure is needed at all. I built
that compiler, the hardware that executes its stream, a banked scratchpad with
cycle-exact conflict counters, and a page-table scheme to compare against ---
and I measured both on identical RTL. Result: 73.2\,\% less off-chip read
traffic than always-reload, and $42\times$ fewer bank conflicts than the paged
scheme on the same layer. I also report where paged wins, and why.
\end{keybox}

\tableofcontents
\vspace{0.4em}
\hrule

% =====================================================================
\section{Research problem}

\subsection{The problem in one sentence}

\begin{keybox}{Research problem}
For a systolic-array DNN accelerator, \textbf{how should the limited on-chip
scratchpad be managed} --- who decides which tensor tiles occupy which physical
addresses, and when --- so that off-chip DRAM traffic, address-translation
overhead and SRAM bank conflicts are all minimised simultaneously?
\end{keybox}

\subsection{Why it matters --- the quantitative motivation}

A single convolution layer with $C{=}16$ input channels, $K{=}32$ output
channels, a $14\times14$ feature map and $3\times3$ kernels needs
\[
  \underbrace{C\,H\,W}_{\text{inputs}} +
  \underbrace{K\,C\,K_h\,K_w}_{\text{weights}} +
  \underbrace{K\,O_h\,O_w}_{\text{outputs}}
  = 3136 + 4608 + 6272 = 14{,}016 \text{ elements}
\]
which at 4 bytes/element is \textbf{56\,KB against a 16\,KB scratchpad}. The
tensors must be tiled, and the accelerator must continuously decide what stays
resident. That single decision drives three costs:

\begin{enumerate}[leftmargin=1.5em,itemsep=3pt]
  \item \textbf{Off-chip traffic.} A DRAM byte costs roughly two orders of
        magnitude more energy than an on-chip SRAM byte ($\approx$560\,pJ/B
        versus a few pJ). Redundant re-fetches dominate the energy budget.
  \item \textbf{Address-translation overhead.} Any runtime structure mapping
        logical tile addresses to physical scratchpad locations charges a
        lookup on \emph{every} access, replicated per read port.
  \item \textbf{Bank conflicts.} The scratchpad is banked so several ports can
        read per cycle. When two ports hit the same bank they serialise and
        the array stalls.
\end{enumerate}

\subsection{The sub-questions this module answers}

\begin{enumerate}[leftmargin=1.5em,itemsep=2pt]
  \item Can a compiler eliminate the runtime tag/translation structure entirely,
        and what does that cost instead?
  \item How much off-chip read traffic does compile-time delta fetching avoid,
        across layer shapes, scratchpad capacities and array sizes?
  \item Static versus dynamic management --- what is the real trade in area,
        clock frequency, cycles, traffic and conflicts?
  \item Are bank conflicts a random phenomenon (in which case nothing can be
        done) or a deterministic function of the access pattern (in which case
        compile-time placement can control them)?
\end{enumerate}

% =====================================================================
\section{Research gap}

\begin{keybox}{Research gap}
No existing DNN-accelerator simulator models off-chip traffic,
address-translation cost and bank conflicts \emph{cycle-accurately in the same
framework}, and none of them supports a compiler-assisted tagless management
scheme at all. Every published comparison of memory-management schemes for
accelerators is therefore either analytical or restricted to cache/page-table
variants.
\end{keybox}

\subsection{What the existing tools do and do not do}

\begin{center}
\small
\begin{tabular}{@{}lccc@{}}
\toprule
\textbf{Simulator} & \textbf{Off-chip traffic} & \textbf{Bank conflicts} &
\textbf{Tagless / compiler-managed} \\
\midrule
SCALE-Sim            & analytical    & no  & no \\
Timeloop / Accelergy & analytical    & no  & no \\
CiMLoop              & analytical    & no  & no \\
STONNE               & partly cycle-level & no  & no \\
ONNXim               & cycle-level   & no  & no \\
\midrule
\textbf{This module} & \textbf{measured (RTL + AXI counters)} &
\textbf{cycle-exact counters} & \textbf{yes, implemented in RTL} \\
\bottomrule
\end{tabular}
\end{center}

\subsection{The three specific gaps}

\begin{enumerate}[leftmargin=1.5em,itemsep=3pt]
  \item \textbf{Gap in modelling fidelity.} Bank conflicts are almost always
        either ignored or approximated with a closed-form birthday-problem
        bound. That bound assumes \emph{uniformly random} addresses; real DNN
        access patterns are highly structured, so the bound is wrong in both
        directions --- consecutive access is conflict-free (far better than the
        bound) and stride-$B$ access is fully serialised (far worse). Nobody
        had measured which one real tile footprints produce.
  \item \textbf{Gap in the design space.} The literature compares caches,
        double buffers and scratchpads with \emph{runtime} allocation. The
        compiler-assisted static tagless point --- no tag array, no page table,
        placement fixed at compile time --- is discussed conceptually but has
        not been implemented in RTL and measured against a dynamic scheme on
        identical hardware.
  \item \textbf{Gap in evidence.} Scheme comparisons in accelerator papers
        typically quote normalised factors from prior work. There was no
        artefact where the same layer, the same array, the same clock and the
        same counters produce both schemes' numbers.
\end{enumerate}

\begin{warnbox}{Say this if asked ``isn't delta encoding old?''}
Yes --- delta encoding itself is decades old. The gap is not the idea of a
delta; it is (a) applying it to \emph{scratchpad layout} rather than to data
values, (b) using it to remove the runtime translation structure completely,
and (c) measuring the consequence cycle-accurately against a page table. It is
the third that no prior simulator provides.
\end{warnbox}

% =====================================================================
\section{Novelty}

\begin{keybox}{The three novel claims}
\textbf{N1 --- Delta-of-layouts, not delta-of-data.} The compiler computes
the difference between consecutive \emph{scratchpad layouts} (which tile sits
at which physical address in each phase) and emits it as a five-opcode stream.
The unit of the delta is a tile placement, not a data value.\\[4pt]
\textbf{N2 --- Tagless by construction.} Because a DNN's loop nest is fully
known before execution, there is nothing to \emph{discover} at runtime, so
there is nothing for a tag or a page table to record. The tag array,
comparators and replacement policy are removed, not optimised.\\[4pt]
\textbf{N3 --- Measured, not modelled.} Both schemes run on the same RTL, and
every reported number comes from a hardware counter, an AXI responder or a
Vivado synthesis report --- including the numbers that do \emph{not} favour my
scheme.
\end{keybox}

\subsection{Supporting novelty}

\begin{itemize}[leftmargin=1.5em,itemsep=3pt]
  \item \textbf{The \op{zero} opcode.} Padded convolutions have receptive
        fields extending past the tensor border. Those bytes are implicit
        zeros: they must exist on-chip but must never be fetched from DRAM.
        Giving them their own opcode --- rather than charging them as loads ---
        reduced the reported load volume from 235{,}008\,B to 157{,}888\,B and
        made the measurement honest rather than flattering.
  \item \textbf{Bank conflicts made a first-class, controllable quantity.}
        The banked scratchpad exposes cycle-exact conflict and stall-port-cycle
        counters, and the measured sweep proves conflicts are a deterministic
        function of the address pattern --- which is exactly what compile-time
        placement controls. This converts ``bank conflicts'' from an assumed
        overhead into a design variable.
  \item \textbf{Three-way independent agreement} (compiler / software FSM
        model / SystemVerilog RTL) as the correctness argument, rather than a
        single golden vector.
\end{itemize}

% =====================================================================
\section{My contribution}

Concretely, what I built and what I measured.

\subsection{Compiler side}

\begin{itemize}[leftmargin=1.5em,itemsep=3pt]
  \item \file{stamp\_compiler.py} (985 lines) --- decomposes a conv layer into
        \emph{phases}, allocates a \emph{stamp} (a complete scratchpad layout)
        per phase, computes the \emph{delta} between consecutive stamps, and
        emits three artefacts: \file{delta\_ops.hex} (the hardware op stream),
        \file{stamp\_phase\_table.svh} (phase base addresses) and
        \file{stamp\_metadata.json} (the machine-readable schedule).
  \item Classification rules based on tile \emph{identity} rather than shape:
        weights are reusable only when the output-channel slice matches;
        input reuse is computed from receptive-field overlap; a simultaneous
        row-and-column shift conservatively falls back to a full reload, so
        \textbf{measured reuse is a lower bound, never an overstatement}.
  \item Padding clipping: the region is intersected with the real tensor
        extent and split into \op{load} bytes and \op{zero} bytes.
  \item \file{stamp\_analyzer.py} --- reuse and bandwidth analysis over the
        emitted stream.
\end{itemize}

\subsection{Hardware side (RTL)}

\begin{itemize}[leftmargin=1.5em,itemsep=3pt]
  \item \file{stamp\_based\_memory\_controller.sv} --- a nine-state FSM that
        reads the compiler's metadata word by word and executes it: issues AXI
        bursts for \op{load}, on-chip read/write pairs for \op{move}, nothing
        for \op{keep} and \op{alloc}, and write-only zero fills for \op{zero}.
        Exposes \file{stats\_loads}, \file{stats\_moves}, \file{stats\_keeps},
        \file{stats\_bytes\_loaded}, \file{stats\_bytes\_moved} and
        \file{stats\_bad\_ops}.
  \item \file{scratchpad\_ram.sv} --- dual-mode SRAM. \file{NUM\_BANKS==1} is a
        flat, conflict-free baseline; \file{NUM\_BANKS}$\,\ge 2$ is low-order
        interleaved with per-cycle conflict arbitration (lowest-indexed port
        wins, losers get \file{rd\_valid=0} and retry) and three monitoring
        outputs: \file{bank\_conflict\_detected},
        \file{stats\_bank\_conflicts},
        \file{stats\_bank\_conflict\_stall\_cycles}.
  \item \file{paged\_memory\_backend.sv} + \file{paged\_table.v} --- the
        dynamic comparison scheme: per-port page tables, VPN$\to$PPN
        translation, hit/miss counters.
  \item \file{mem\_backend\_wrap.sv} --- selects STAMP or PAGED by parameter
        and propagates all counters to the top level.
  \item \textbf{Plumbing fixes I made so the comparison is real:}
        \file{NUM\_BANKS} is now a genuine top-level parameter (it was
        hard-defaulted to 4 and never reached the simulator); the bank counters
        were dangling at \file{u\_mem} and are now top ports; and external
        \file{spad\_dbg\_*} read ports were added so the memory side can be
        exercised independently of the array.
  \item \textbf{PAGED page-table write bug (found and fixed).} The per-port PT
        write was gated on \verb|p == 0| \emph{and} required \file{rd\_addr[0]}
        to hold the target VPN, so ports 1--3 could never hit --- the paged
        scheme was structurally crippled. Writes now broadcast to all port
        tables. \textbf{The whole comparison was re-run against this corrected
        baseline}, i.e. against a paged scheme working as well as it can.
\end{itemize}

\subsection{Model and evaluation side}

\begin{itemize}[leftmargin=1.5em,itemsep=3pt]
  \item \file{pysim/stamp\_ref.py} --- an independently written cycle-accurate
        model of the controller FSM, used both as a cross-check and as the
        engine for large sweeps where RTL simulation is too slow.
  \item \file{scripts/sweep\_memory\_schemes.py} --- 6 layer shapes $\times$ 4
        scratchpad capacities $\times$ 2 array sizes; each grid point is a real
        compiler run, not an interpolation. 34 of 48 points fit; the other 14
        are recorded as capacity failures rather than silently dropped.
  \item \file{tb/golden/test\_scheme\_divergence.py} --- the STAMP-vs-PAGED
        head-to-head on real layers with real tile footprints.
  \item \file{tb/golden/test\_bank\_sweep.py} --- the four-pattern $\times$
        five-bank-count conflict sweep.
  \item \file{scripts/run\_full\_eval.py} $\to$ \file{exp4\_memory\_management}
        and \file{exp4b\_stamp\_vs\_paged} --- turn the captured artefacts into
        the report figures.
  \item \textbf{De-hardcoding.} Following the supervisor's ``don't hardcode
        anything'' instruction I removed the previous four-scheme comparison
        entirely (see \S\ref{sec:superseded}) and replaced it with measured
        static-vs-dynamic. I also removed the invented constants in the bank
        conflict experiment (\file{WORKLOAD\_MEM\_INTENSITY},
        \file{DATAFLOW\_PORT\_ACTIVITY}, and the clamps that pinned every conv
        layer to exactly 0.18) and replaced them with quantities derived from
        each layer's own geometry.
\end{itemize}

% =====================================================================
\section{Artefacts}

\subsection{Source artefacts (what I wrote)}

\begin{center}
\small
\begin{longtable}{@{}p{7.2cm}p{8.4cm}@{}}
\toprule
\textbf{File} & \textbf{Role} \\
\midrule
\endhead
\file{static\_hash\_and\_tagless\_memory/stamp\_compiler.py} & Phases, stamps, deltas; emits hex + phase table + metadata \\
\file{.../stamp\_analyzer.py} & Reuse and bandwidth analysis \\
\file{.../stamp\_based\_memory\_controller.sv} & 9-state delta-op sequencer \\
\file{.../tb\_stamp\_based\_memory\_controller.sv} & 13-test directed testbench (48 assertions) \\
\file{.../paged\_table.v}, \file{tb\_paged\_table.v} & Page table + its testbench \\
\file{.../run\_vivado\_sim.sh} & One-command reproduction of the RTL evidence \\
\file{.../synth/synth\_area.tcl} & Vivado out-of-context synthesis driver \\
\file{sim\_framework/rtl/memory/scratchpad\_ram.sv} & Banked SRAM + conflict counters \\
\file{.../rtl/memory/paged\_memory\_backend.sv} & Dynamic (page-table) backend \\
\file{.../rtl/memory/stamp\_memory\_backend.sv} & Static (STAMP) backend \\
\file{.../rtl/memory/mem\_backend\_wrap.sv} & Scheme selector + counter propagation \\
\file{sim\_framework/pysim/stamp\_ref.py} & Cycle-accurate FSM software model \\
\file{sim\_framework/tb/unit/test\_stamp\_controller.py} & 27 software tests \\
\file{sim\_framework/tb/unit/test\_bank\_conflict.py} & 17 banking tests \\
\file{sim\_framework/tb/golden/test\_scheme\_divergence.py} & STAMP vs PAGED on real layers \\
\file{sim\_framework/tb/golden/test\_bank\_sweep.py} & Bank-conflict pattern sweep \\
\file{sim\_framework/scripts/sweep\_memory\_schemes.py} & 34-point compiler sweep \\
\file{sim\_framework/scripts/run\_full\_eval.py} & \file{exp3}, \file{exp4}, \file{exp4b} figure generation \\
\file{.../docs/memory\_module\_theory.pdf} & Full theory + architecture document \\
\bottomrule
\end{longtable}
\end{center}

\subsection{Generated artefacts (the evidence)}

\begin{center}
\small
\begin{longtable}{@{}p{8.0cm}p{7.6cm}@{}}
\toprule
\textbf{Artefact} & \textbf{What it holds} \\
\midrule
\endhead
\file{static\_hash\_and\_tagless\_memory/delta\_ops.hex} & The compiler's op stream, read by RTL via \file{\$readmemh} \\
\file{.../stamp\_phase\_table.svh} & Phase base addresses, \file{`include}d by RTL \\
\file{.../stamp\_metadata.json} & Machine-readable schedule (drives the accounting figures) \\
\file{.../stamp\_analysis\_summary.txt} & Reuse / bandwidth summary of the standalone run \\
\file{.../synth/work/out/area\_results.csv} & \textbf{Vivado OOC synthesis}: LUTs, FFs, WNS, Fmax for 11 configurations \\
\file{sim\_framework/results/golden\_check/raw/divergence\_*.json} & \textbf{Per-run RTL counters} for STAMP and PAGED (4 files) \\
\file{.../results/golden\_check/raw/bank\_sweep\_b\{1,2,4,8,16\}.json} & \textbf{Cycle-exact} conflict/stall counters per pattern \\
\file{.../results/exp4\_memory\_management/sweep\_stamp.csv} & 34 real compiler runs across the design grid \\
\file{.../results/exp4\_memory\_management/capacity\_limits.csv} & The 14 grid points that do not fit, with the reason \\
\file{.../results/exp4\_memory\_management/area\_measured.csv} & Copy of the synthesis table used by the figures \\
\file{.../results/exp4b\_stamp\_vs\_paged/stamp\_vs\_paged.csv} & The head-to-head comparison table \\
\file{.../results/exp4b\_stamp\_vs\_paged/bank\_conflict\_sweep.csv} & The 20-row bank sweep table \\
\file{.../results/exp4b\_stamp\_vs\_paged/stamp\_delta\_accounting.csv} & Load / move / zero byte accounting \\
\file{.../results/vivado/stamp\_controller\_xsim.log} & The xsim console log (48 assertions, 0 failures) \\
\file{.../results/golden\_check/MEMBER2\_MEMORY\_GOLDEN\_REPORT.md} & The written golden-check report with honesty labels \\
\bottomrule
\end{longtable}
\end{center}

\subsection{Figures produced}

\file{results/exp4\_memory\_management/figures/}: \file{area\_static\_vs\_dynamic.png},
\file{area\_scaling.png}, \file{banking\_area\_cost.png},
\file{offchip\_reduction\_sweep.png}, \file{metadata\_overhead.png}.

\file{results/exp4b\_stamp\_vs\_paged/figures/}: \file{axi\_beats.png},
\file{axi\_ar\_requests.png}, \file{bank\_conflicts.png},
\file{stall\_port\_cycles.png}, \file{cycle\_breakdown.png},
\file{bank\_stall\_port\_cycles.png}, \file{bank\_reads\_served.png},
\file{stamp\_vs\_naive\_traffic.png}, \file{delta\_op\_mix.png}.

% =====================================================================
\section{Methodology}

\subsection{Overall flow}

\begin{center}
\small
\begin{tabular}{@{}rp{13.5cm}@{}}
\toprule
\textbf{Step} & \textbf{What happens} \\
\midrule
1 & \textbf{Compile.} \file{stamp\_compiler.py} takes a layer geometry and a
    scratchpad size, builds the phase list, allocates a stamp per phase, and
    computes deltas. Emits \file{delta\_ops.hex}, \file{stamp\_phase\_table.svh},
    \file{stamp\_metadata.json}. \\
2 & \textbf{Model.} \file{pysim/stamp\_ref.py} replays that op stream through a
    cycle-accurate model of the controller FSM and returns cycles, bytes loaded,
    bytes moved, bytes zeroed, AXI bursts. \\
3 & \textbf{Simulate (RTL, xsim).} The SystemVerilog controller reads the
    \emph{same} hex file via \file{\$readmemh} and executes all 127 phase
    transitions / 572 operations under the directed testbench. \\
4 & \textbf{Simulate (RTL, Verilator + cocotb).} The full accelerator top runs a
    real layer end-to-end, once with \file{MEMORY=STAMP} and once with
    \file{MEMORY=PAGED}, at \file{NUM\_BANKS} $\in \{1,2,4,8,16\}$, and dumps
    every hardware counter to JSON. Correctness is checked against a
    TensorFlow golden reference. \\
5 & \textbf{Synthesise.} Vivado out-of-context synthesis of the stamp
    controller, the paged manager and the banked scratchpad, swept over each
    one's own capacity knob, yields LUT/FF/Fmax. \\
6 & \textbf{Sweep.} \file{sweep\_memory\_schemes.py} runs steps 1--2 over the
    full layer $\times$ capacity $\times$ array grid. \\
7 & \textbf{Report.} \file{run\_full\_eval.py} reads only the artefacts from
    steps 4--6 and produces the CSVs and figures. It contains no performance
    constants. \\
\bottomrule
\end{tabular}
\end{center}

\subsection{Experimental configuration}

\begin{center}
\small
\begin{tabular}{@{}ll@{}}
\toprule
\textbf{Parameter} & \textbf{Value} \\
\midrule
Layers (head-to-head)   & \file{tiny\_cnn} layer 0, \file{mnist\_cnn} layer 0 \\
Dataflow / layout       & Output-stationary, $8\times8$ array, CHANNEL\_MAJOR \\
Scratchpad read ports   & 4 \\
Bank counts swept       & 1, 2, 4, 8, 16 \\
Page size (PAGED)       & $2^{10}$ words = 4\,KB \\
Metadata depth (STAMP)  & 256 / 512 / 1024 entries (1024 used for verification) \\
Standalone compiler run & 16\,KB scratchpad, 4\,B elements, 128 phases \\
Sweep grid              & 6 layer shapes $\times$ \{8,16,32,64\}\,KB $\times$ \{$4\times4$, $8\times8$\} \\
Tools                   & Verilator 5.050 + cocotb 2.0.1, Vivado xsim 2025.2, TF 2.21 \\
Correctness tolerance   & 5\,\% of full scale (achieved: 0.009--0.013\,\%) \\
\bottomrule
\end{tabular}
\end{center}

\subsection{The two conflict models, and which is authoritative}

\begin{warnbox}{Be precise about this in the viva}
There are two bank-conflict numbers in the report and they are not the same
thing:
\begin{itemize}[leftmargin=1.3em,itemsep=2pt]
  \item \textbf{Experiment 3} uses the analytical birthday-problem bound,
        $P_{\text{conflict}} = 1 - \prod_{k=0}^{P-1}\frac{B-k}{B}$, scaled by a
        per-layer memory-active fraction derived from the layer's own MAC count
        and byte demand. It is an \emph{estimate}, applied across 14 workloads
        where RTL simulation of every layer would be prohibitive.
  \item \textbf{Experiment 4b} uses \emph{cycle-exact hardware counters} from
        the banked scratchpad. It is the authority.
\end{itemize}
The bound assumes uniform random addresses; real patterns are structured, so
the measured values differ from the bound --- and that difference is itself one
of my findings.
\end{warnbox}

% =====================================================================
\section{Alternative memory management schemes --- the comparison}
\label{sec:compare}

\subsection{What is compared against what}

Three comparison axes, all measured:

\begin{center}
\small
\begin{tabular}{@{}p{2.6cm}p{5.2cm}p{7.4cm}@{}}
\toprule
\textbf{Axis} & \textbf{Alternatives} & \textbf{Where measured} \\
\midrule
Management policy & \textbf{STAMP} (static, tagless, compiler-assisted)
                    vs \textbf{PAGED} (dynamic, hardware page table)
                  & \file{exp4b} --- RTL counters; \file{exp4} --- Vivado area/Fmax \\
Fetch policy      & \textbf{Delta fetch} vs \textbf{always-load}
                    (a tag-based scheme that re-reads every input and weight
                    tile every phase)
                  & \file{exp4} sweep + \file{stamp\_delta\_accounting.csv} \\
Memory organisation & \textbf{Banking degree} $B \in \{1,2,4,8,16\}$
                    ($B{=}1$ = flat, conflict-free baseline)
                  & \file{bank\_conflict\_sweep.csv} + \file{area\_results.csv} \\
\bottomrule
\end{tabular}
\end{center}

\subsection{The two schemes, mechanically}

\begin{center}
\small
\begin{tabular}{@{}p{3.5cm}p{5.8cm}p{5.9cm}@{}}
\toprule
 & \textbf{PAGED (dynamic)} & \textbf{STAMP (static)} \\
\midrule
Who decides placement & Hardware, at runtime & Compiler, ahead of time \\
Runtime structure & Page table per port; VPN $\to$ PPN + valid bit & None --- direct addressing \\
Lookup cost & Every access, every port & Zero \\
Fetch granularity & Whole page (4\,KB) on a fault & Exactly the new bytes \\
Bookkeeping storage & Table $=$ NUM\_PAGES $\times$ (PPN width $+$ 1) bits & Metadata RAM $=$ 128 bits per delta op \\
Bank residency & Unpredictable to the compiler & Chosen by the compiler \\
Latency determinism & Non-deterministic (faults) & Deterministic \\
\bottomrule
\end{tabular}
\end{center}

\subsection{The five delta operations (STAMP)}

\begin{center}
\small
\begin{tabular}{@{}llll@{}}
\toprule
\textbf{Op} & \textbf{Opcode} & \textbf{Condition} & \textbf{Cost} \\
\midrule
\op{keep}  & 0 & identical data, identical address & 0 cycles, 0 bytes \\
\op{move}  & 1 & identical data, different address & 2 cycles/word (rd+wr) \\
\op{load}  & 2 & data not resident on-chip & \textbf{DRAM burst} \\
\op{alloc} & 3 & output slot, written by compute & 0 cycles, 0 bytes \\
\op{zero}  & 4 & zero-padding halo & 1 cycle/word (write only) \\
\bottomrule
\end{tabular}
\end{center}

\noindent Only \op{load} costs off-chip bandwidth. The saving is
$\sigma = R/(L+R)$ with $R = M+K$ the bytes resolved on-chip.

\subsection{Superseded: the four-scheme comparison}
\label{sec:superseded}

\begin{warnbox}{If an examiner has seen the older charts}
An earlier version of Experiment 4 compared four schemes --- Static, Double
Buffer, Unified Buffer, Cache-Based. \textbf{Every number in it came from a
hardcoded dictionary} (\file{MEMSCHEME\_MODEL}: area 1.18 / cycles 0.76, etc.)
and invented formulas of the form $0.90 - 0.28\,m_f$. Two problems: nothing was
measured, and those three scheme names appear nowhere in
\file{evaluation\_requirements.pdf}, which asks for ``tradeoffs between static
and dynamic memory management schemes''. On the supervisor's ``don't hardcode
anything'' instruction (Aug 2026) the whole experiment was deleted and replaced
with the measured static-vs-dynamic comparison. The old outputs are retained,
clearly labelled, under \file{results/\{cloud,edge\}/\_superseded\_hardcoded\_exp4/}
so the change is auditable. \textbf{Do not present those figures.}
\end{warnbox}

% =====================================================================
\section{The compared values --- what they are, how they are captured}

\subsection{Metric dictionary}

\begin{center}
\small
\begin{longtable}{@{}p{4.3cm}p{5.0cm}p{6.2cm}@{}}
\toprule
\textbf{Value} & \textbf{What it means} & \textbf{Source / how captured} \\
\midrule
\endhead
\file{stats\_bytes\_loaded} & Bytes actually fetched from DRAM by the memory backend & Hardware counter inside \file{stamp\_memory\_backend}, incremented per accepted AXI beat \\
\file{axi\_ar\_requests} & AXI read-address handshakes issued & Counted by the cocotb AXI responder model \\
\file{axi\_beats} & AXI data beats returned & Counted by the cocotb AXI responder model \\
\file{stats\_bank\_conflicts} & Cycles in which $\ge 1$ port was denied its bank & Counter in \file{scratchpad\_ram.sv}, increments when \file{bank\_conflict\_detected} \\
\file{stats\_bank\_conflict\_stall\_cycles} & Total \emph{port}-cycles stalled (sum over stalled ports) & Same block; adds the stalled-port count each conflict cycle \\
\file{compute\_cycles} & Cycles the top FSM spends in compute for the whole layer & Free-running RTL cycle counter, sampled at layer end \\
\file{program\_cycles} & Cycles spent writing the metadata / page table & Same counter, sampled around the programming phase \\
\file{replay\_cycles} & Cycles of the 4-port footprint replay & Same counter, sampled around the replay \\
\file{stats\_loads\_or\_hits}, \file{stats\_moves\_or\_misses}, \file{stats\_keeps} & STAMP: load/move/keep op counts. PAGED: page hits/misses & Backend counters; the wrapper reuses the same ports for both schemes \\
\file{new\_words} / \file{reused\_words} & Word-touches classified new vs already resident & Host-side bookkeeping in the testbench, cross-checked against the HW counter \\
\file{luts}, \file{ffs}, \file{wns\_ns}, \file{fmax\_mhz} & Post-synthesis resource and timing & Vivado out-of-context synthesis reports \\
\file{reduction\_pct} & $1 - L/\text{naive}$, off-chip read traffic avoided & Compiler + FSM model over the sweep grid \\
\file{metadata\_bytes} & On-chip RAM the delta stream needs & $16 \times$ (total delta ops) \\
\bottomrule
\end{longtable}
\end{center}

\subsection{The capture process, step by step}

\paragraph{(a) RTL counters for the head-to-head (\file{divergence\_*.json}).}
\begin{enumerate}[leftmargin=1.5em,itemsep=2pt]
  \item \file{tb/golden/run\_golden.py} builds the design with Verilator,
        passing \file{-GMEMORY=\{0,1\}} (STAMP/PAGED) and
        \file{-GNUM\_BANKS=\textit{n}} as parameter overrides.
  \item \file{test\_scheme\_divergence.py} walks the layer's OS tile sequence.
        For \textbf{STAMP} it tracks which DRAM words are already resident and
        programs the metadata table with \op{load} ops for contiguous runs of
        new words and \op{keep} ops for contiguous runs of reused words; the
        hardware engine then executes them during the stamp phase, issuing real
        AXI bursts. For \textbf{PAGED} it replays the tile footprint's DRAM
        addresses through the translation hardware, so unmapped pages MISS and
        mapped pages HIT, then maps the faulted pages with \file{pt\_write}.
  \item After each tile, the tile footprint is replayed on all four external
        \file{spad\_dbg\_*} read ports, so the banked arbitration sees real DNN
        address patterns.
  \item At the end, every counter is read out and written to
        \file{results/golden\_check/raw/divergence\_\{layer\}\_\{scheme\}\_b\{n\}.json},
        together with the full output tensor and a pass/fail verdict against
        the TensorFlow reference.
\end{enumerate}

\paragraph{(b) Bank sweep (\file{bank\_sweep\_b*.json}).}
\file{test\_bank\_sweep.py} never starts the DUT FSM --- it measures the memory
side alone. It drives the four external read ports with a fixed address pattern
for a 200-cycle window and reads the hardware counters before and after:
\begin{center}\small
\begin{tabular}{@{}ll@{}}
\toprule
\textbf{Pattern} & \textbf{Port $p$ reads address} \\
\midrule
\file{same\_addr}   & 0 (all ports) \\
\file{consecutive}  & $p$ \\
\file{stride4}      & $4p$ \\
\file{stride8}      & $8p$ \\
\bottomrule
\end{tabular}
\end{center}
Repeated for $B \in \{1,2,4,8,16\}$ (five separate builds).

\paragraph{(c) Area and Fmax (\file{area\_results.csv}).}
\file{synth/synth\_area.tcl} runs Vivado out-of-context synthesis on each module
in isolation (no I/O buffer insertion, so the numbers are the logic's own cost),
sweeping the stamp controller over \file{METADATA\_DEPTH} $\in$ \{256,512,1024\},
the paged manager over \file{NUM\_PAGES} $\in$ \{128,256,512\} and the
scratchpad over \file{NUM\_BANKS} $\in$ \{1,2,4,8,16\}. The utilisation and
timing reports are parsed into one CSV row per configuration.

\paragraph{(d) Compiler sweep (\file{sweep\_stamp.csv}).}
\file{scripts/sweep\_memory\_schemes.py} instantiates \file{StampCompiler} at
each grid point, calls \file{create\_conv\_phases} $\to$
\file{allocate\_stamps} $\to$ \file{compute\_deltas}, then replays the emitted
ops through \file{pysim.stamp\_ref.run\_stream}. If \file{allocate\_stamps}
raises (the phase does not fit), the point is recorded in
\file{capacity\_limits.csv} instead of being dropped.

\paragraph{(e) Standalone compiler-to-RTL evidence.}
\file{run\_vivado\_sim.sh} runs the compiler, runs the analyzer, then compiles
and simulates \file{tb\_stamp\_based\_memory\_controller.sv} in xsim against the
generated \file{delta\_ops.hex}, and copies the console log into
\file{results/vivado/}.

\subsection{Reproduction commands}

\begin{center}
\small
\begin{tabular}{@{}p{5.3cm}p{10.2cm}@{}}
\toprule
\textbf{To regenerate} & \textbf{Command} \\
\midrule
Compiler artefacts + xsim log & \file{./static\_hash\_and\_tagless\_memory/run\_vivado\_sim.sh} \\
Area / Fmax CSV & \file{vivado -mode batch -source synth/synth\_area.tcl} \\
Compiler sweep CSVs & \file{PYTHONPATH=. python scripts/sweep\_memory\_schemes.py} \\
RTL head-to-head + bank sweep & \file{PYTHONPATH=. python tb/golden/run\_golden.py --memory \{STAMP,PAGED\} --num-banks \textit{n} --layer \ldots} \\
Software test suites & \file{PYTHONPATH=. python -m pytest tb/unit/test\_stamp\_controller.py tb/unit/test\_bank\_conflict.py -v} \\
All figures & \file{PYTHONPATH=. python scripts/run\_full\_eval.py} \\
\bottomrule
\end{tabular}
\end{center}

% =====================================================================
\section{Comparison results}

\subsection{R1 --- STAMP vs PAGED, measured RTL counters}

Whole layer, OS $8\times8$, CHANNEL\_MAJOR, 4 banks. Source:
\file{results/exp4b\_stamp\_vs\_paged/stamp\_vs\_paged.csv}.

\begin{center}
\small
\begin{tabular}{@{}lrrrr@{}}
\toprule
\textbf{Metric} & \textbf{tiny STAMP} & \textbf{tiny PAGED} &
                  \textbf{mnist STAMP} & \textbf{mnist PAGED} \\
\midrule
Correctness (max rel.\ err)   & PASS 0.013\,\% & PASS 0.013\,\% & PASS 0.009\,\% & PASS 0.009\,\% \\
Bank-conflict events          & \textbf{2}   & 84    & \textbf{101}  & 232 \\
Stalled port-cycles           & \textbf{4}   & 240   & \textbf{202}  & 488 \\
Backend off-chip bytes        & 1{,}640 (meas.) & 8{,}192 (derived) & 3{,}448 (meas.) & 8{,}192 (derived) \\
AXI read requests             & 157  & \textbf{138}   & 945   & \textbf{832} \\
AXI beats                     & 2{,}246 & \textbf{1{,}836} & 11{,}470 & \textbf{10{,}608} \\
Compute cycles                & 6{,}316 & 5{,}802 & 38{,}440 & 36{,}608 \\
Metadata / PT programming cyc & 41   & 2     & 527   & 2 \\
Replay cycles                 & 456  & 456   & 2{,}912 & 2{,}912 \\
\op{load} / \op{keep} ops     & 19 / 22 & --- & 113 / 414 & --- \\
Page hits / misses            & ---  & 1{,}470 / 294 & --- & 10{,}506 / 102 \\
New / reused word-touches     & 410 / 1{,}354 & 410 / 1{,}354 & 862 / 9{,}746 & 862 / 9{,}746 \\
\bottomrule
\end{tabular}
\end{center}

\begin{keybox}{Headline}
\textbf{Bank conflicts: $42\times$ fewer} (2 vs 84) on tiny, $2.3\times$ fewer
on mnist. \textbf{Stalled port-cycles: $60\times$ fewer} (4 vs 240) on tiny.
\textbf{Off-chip volume: $5.0\times$ less} on tiny (1{,}640\,B vs 8{,}192\,B)
and $2.4\times$ less on mnist (3{,}448\,B vs 8{,}192\,B). STAMP's delta
mechanism captured \textbf{76.8\,\%} (tiny) and \textbf{91.9\,\%} (mnist) of
word-touches as on-chip reuse. \textbf{But PAGED transfers fewer AXI beats on
these layers} --- and I report that.
\end{keybox}

\subsection{R2 --- Off-chip traffic vs the always-load baseline}

Standalone 128-phase layer, 16\,KB scratchpad. Source:
\file{stamp\_analysis\_summary.txt}, \file{stamp\_delta\_accounting.csv}.

\begin{center}
\small
\begin{tabular}{@{}lrr@{}}
\toprule
\textbf{Quantity} & \textbf{Bytes} & \textbf{Share} \\
\midrule
Always-load baseline (inputs + weights, every phase) & 589{,}824 & 100.0\,\% \\
STAMP scheme, actual DRAM reads                      & 157{,}888 & 26.8\,\% \\
\midrule
Resolved on-chip by \op{move} & 184{,}320 & --- \\
Resolved on-chip by \op{keep} & 165{,}888 & --- \\
Zero-filled padding (never off-chip) & 77{,}120 & --- \\
\midrule
\textbf{Off-chip read traffic reduction} & \multicolumn{2}{r}{\textbf{73.2\,\%}} \\
Bandwidth saving $\sigma = R/(L{+}R)$    & \multicolumn{2}{r}{68.9\,\%} \\
\bottomrule
\end{tabular}
\end{center}

\noindent Op mix: 134 \op{load}, 144 \op{move}, 72 \op{keep}, 127 \op{alloc},
95 \op{zero} $=$ 572 operations over 127 phase transitions.

\subsection{R3 --- The design-space sweep (34 real compiler runs)}

Source: \file{results/exp4\_memory\_management/sweep\_stamp.csv}.

\begin{center}
\small
\begin{tabular}{@{}lrrr@{}}
\toprule
\textbf{Layer shape} & \textbf{Min red.} & \textbf{Mean red.} & \textbf{Max red.} \\
\midrule
early-conv $3\times3$   & 64.2\,\% & 66.1\,\% & 67.5\,\% \\
mid-conv $3\times3$     & 59.8\,\% & 67.6\,\% & 72.9\,\% \\
late-conv $3\times3$    & 55.9\,\% & 62.8\,\% & 66.3\,\% \\
pointwise $1\times1$    & 32.3\,\% & 36.0\,\% & 37.9\,\% \\
wide-spatial            & 54.9\,\% & 61.1\,\% & 67.3\,\% \\
large-kernel $5\times5$ & 73.0\,\% & 73.4\,\% & 73.6\,\% \\
\midrule
\textbf{All 34 points}  & \textbf{32.3\,\%} & \textbf{60.6\,\%} & \textbf{73.6\,\%} \\
\bottomrule
\end{tabular}
\end{center}

\noindent Controller-cycle saving against a 10-cycle/word DRAM baseline ranges
88.0--95.6\,\% (mean 91.4\,\%). 14 of the 48 grid points do not fit and are
recorded with the reason in \file{capacity\_limits.csv}.

\subsection{R4 --- Bank-conflict sensitivity (cycle-exact)}

Stalled port-cycles over a 200-cycle window, 4 read ports. Source:
\file{bank\_conflict\_sweep.csv}.

\begin{center}
\small
\begin{tabular}{@{}lrrrrr@{}}
\toprule
\textbf{Access pattern} & \textbf{1 bank} & \textbf{2} & \textbf{4} & \textbf{8} & \textbf{16} \\
\midrule
All ports $\to$ same address & 0 & 600 & 600 & 600 & 600 \\
Stride 4                     & 0 & 600 & 600 & 400 & \textbf{0} \\
Stride 8                     & 0 & 600 & 600 & 600 & 400 \\
Consecutive                  & 0 & 400 & \textbf{0} & \textbf{0} & \textbf{0} \\
\bottomrule
\end{tabular}
\end{center}

\noindent Reads served in the same window: 796 when conflict-free, 199 when
fully serialised --- a $4\times$ throughput swing driven purely by the address
pattern. With \emph{real} traffic (tiny L0, STAMP, whole layer) the conflict /
stall counts were 438/876 at $B{=}2$, 2/4 at $B{=}4$ and $B{=}8$, and 0/0 at
$B{=}16$.

\subsection{R5 --- Area and frequency (Vivado OOC, ZCU104)}

Source: \file{static\_hash\_and\_tagless\_memory/synth/work/out/area\_results.csv}.

\begin{center}
\small
\begin{tabular}{@{}llrrrr@{}}
\toprule
\textbf{Config} & \textbf{Scheme} & \textbf{LUTs} & \textbf{FFs} & \textbf{WNS (ns)} & \textbf{Fmax (MHz)} \\
\midrule
\file{stamp\_ctrl\_m256}  & STAMP (static)  & 6{,}434  & 22{,}198 & 7.533 & 405.4 \\
\file{stamp\_ctrl\_m512}  & STAMP (static)  & 12{,}574 & 43{,}938 & 7.237 & 361.9 \\
\file{stamp\_ctrl\_m1024} & STAMP (static)  & 24{,}680 & 87{,}468 & 7.512 & 401.9 \\
\file{paged\_mgmt\_p128}  & PAGED (dynamic) & 3{,}074  & 2{,}752  & 8.361 & 610.1 \\
\file{paged\_mgmt\_p256}  & PAGED (dynamic) & 6{,}110  & 5{,}440  & 8.317 & 594.2 \\
\file{paged\_mgmt\_p512}  & PAGED (dynamic) & 12{,}106 & 10{,}816 & 8.121 & 532.2 \\
\midrule
\file{scratchpad\_b1}     & Scratchpad      & 71{,}388 & 8{,}357 & 6.341 & 273.3 \\
\file{scratchpad\_b2}     & Scratchpad      & 38{,}186 & 8{,}421 & 8.834 & 857.6 \\
\file{scratchpad\_b4}     & Scratchpad      & 44{,}880 & 8{,}421 & 8.604 & 716.3 \\
\file{scratchpad\_b8}     & Scratchpad      & 35{,}356 & 8{,}421 & 8.875 & 888.9 \\
\file{scratchpad\_b16}    & Scratchpad      & 37{,}713 & 8{,}421 & 9.001 & 1001.0 \\
\bottomrule
\end{tabular}
\end{center}

\noindent At equal capacity (256 entries) STAMP uses $\approx$1.05$\times$ the
LUTs but $4.1\times$ the registers of PAGED, and clocks 32\,\% slower
(405 vs 594\,MHz). Both scale close to linearly with their capacity knob:
STAMP 28.6\,k $\to$ 56.5\,k $\to$ 112.1\,k LUT+FF over $256\to512\to1024$;
PAGED 5.8\,k $\to$ 11.6\,k $\to$ 22.9\,k over $128\to256\to512$.

% =====================================================================
\section{Why the results came out that way}

\subsection{Why STAMP has $42\times$ fewer bank conflicts}

The stamp compiler bump-allocates scratchpad slots, so a tile's words land in a
\emph{dense, contiguous} region. When four ports replay that footprint they read
stride-1 addresses, which under low-order interleaving map to four
\emph{different} banks --- almost conflict-free at $B{=}4$. The page table
preserves page-relative offsets, so a translated \file{paddr} keeps the original
tensor's row stride; four ports reading four rows of a feature map hit
$a$, $a{+}W$, $a{+}2W$, $a{+}3W$, and whenever $W \bmod B = 0$ all four land in
the same bank. That is a genuine consequence of each scheme's placement policy,
not an artefact. The bank sweep (R4) is the controlled proof: the same hardware
gives 0 stalls for consecutive and 600 for stride-4 at $B{=}4$.

\subsection{Why PAGED transfers fewer AXI beats (the result against me)}

A page fault pulls a whole 4\,KB page. For these small test layers the working
set is small enough that a handful of coarse page fetches happens to cover it
efficiently, while STAMP issues many small delta bursts and pays more AXI
request overhead (157 vs 138 AR handshakes). \textbf{My claim is not ``fewer
bytes always''} --- it is fewer conflicts, no translation hardware, and
deterministic placement. As layers grow, a page fault pulls far more data than
the corresponding delta and the byte comparison inverts; the 73.2\,\% reduction
against always-load in R2 is the regime-independent statement.

\subsection{Why STAMP burns more compute cycles (6{,}316 vs 5{,}802)}

Two reasons, and only one is real. (i) STAMP's real delta-load AXI time is
\emph{inside} its compute cycle count ($+514$ cycles on tiny), whereas the paged
backend has no DRAM path in RTL, so its implied page-fetch time is not charged
anywhere. (ii) STAMP genuinely pays a metadata-programming cost (41 vs 2 cycles
on tiny; 527 vs 2 on mnist). \textbf{The cycle comparison is asymmetric in
PAGED's favour and I state so.}

\subsection{Why pointwise $1\times1$ shows the least reduction (32--38\,\%)}

Reuse in this scheme comes from receptive-field overlap between consecutive
phases. A $1\times1$ kernel has no overlap --- the receptive field is a single
pixel, so sliding the window produces an entirely new footprint. Only weight
reuse remains. The scheme's benefit is therefore proportional to kernel extent,
which is exactly why large-kernel $5\times5$ shows the most (73.0--73.6\,\%).
This is a property of the workload, not a weakness of the implementation.

\subsection{Why the flat scratchpad ($B{=}1$) costs the most LUTs and the worst Fmax}

Counter-intuitive but correct: a flat array serving 4 concurrent read ports must
be implemented as replicated/multiplexed distributed RAM, so it burns 71\,k LUTs
and closes at only 273\,MHz. Splitting it into banks lets each bank serve one
port through a narrow crossbar --- 35--45\,k LUTs and 716--1001\,MHz.
\textbf{Banking is not purely a cost paid for parallelism; on this device it is
also cheaper and faster.} The cost of banking is conflicts, not area.

\subsection{Why measured conflict counts are far below the birthday bound}

For $P{=}4$, $B{=}4$ the analytical bound predicts $P_{\text{conflict}} =
90.6\,\%$ of cycles. The measured STAMP run shows 2 conflict events across the
whole layer. The bound assumes uniformly random addresses; real footprints are
consecutive, and consecutive addresses map to distinct banks \emph{by
construction}. The gap between 90.6\,\% and $\approx$0 is precisely the value of
compile-time placement --- and it is why Experiment 3's analytical numbers must
be labelled as an estimate and Experiment 4b's as the authority.

% =====================================================================
\section{Limitations}

\begin{enumerate}[leftmargin=1.6em,itemsep=5pt]
  \item \textbf{Metadata RAM is the scaling limit.} Removing the tag array does
        not make bookkeeping free --- it relocates it to a compile-time
        instruction stream, $\text{bits} = 128 \times \sum_\phi
        |\Delta_{\phi-1\to\phi}|$. For the measured 128-phase layer, 572 ops
        $\times$ 16\,B $=$ 9{,}152\,B for a 16\,KB scratchpad --- about
        \textbf{56\,\% of the scratchpad size}. Across the sweep grid it ranges
        0.7\,KB to 51\,KB. The default \file{METADATA\_DEPTH} of 256 entries
        holds only $\approx$45\,\% of that layer's stream, so verification uses
        1024. \textbf{This is the honest headline trade-off.}
  \item \textbf{Asymmetric off-chip metric.} STAMP's bytes are measured on the
        real AXI port; PAGED's 8{,}192\,B is \emph{derived} (2 measured page
        faults $\times$ 4\,KB) because the paged backend has no DRAM path in
        RTL. The volume comparison is real; the cycle comparison favours PAGED.
  \item \textbf{v1 wiring limit --- conflicts do not yet back-pressure compute.}
        All conflicts, stalls, hits and misses are memory-side events. The array
        currently reads via the \file{ext\_*} ports rather than through the
        scratchpad, so a stall is counted but does not throttle the pipeline.
        Full compute-coupling is future work.
  \item \textbf{Replay stands in for the array's reads.} The 4-port replays use
        real tile-footprint addresses and real RTL arbitration, but they model
        the consumers rather than being them (consequence of item 3).
  \item \textbf{Delta model handles only pure row or pure column shifts.} A
        simultaneous row-and-column shift falls back to a full reload, so
        reported reuse is a \textbf{lower bound}.
  \item \textbf{No eviction or replacement policy in PAGED (v1).} Pages map once
        and stay. Fine at these footprints ($\ll$ scratchpad); it matters for
        larger networks and slightly flatters PAGED's miss count.
  \item \textbf{Address-unit inconsistency (documented, open).} The stamp
        datapath is byte-addressed (dst \verb|>>2| before the scratchpad); the
        paged datapath passes \file{paddr} low bits to the word-addressed
        scratchpad with no shift. The divergence experiment therefore runs PAGED
        word-addressed with \file{PAGE\_SIZE\_BITS=10} for self-consistency.
  \item \textbf{Two op-stream generators still exist.} The standalone
        compiler$\to$RTL path is fully closed (\file{delta\_ops.hex} read by
        \file{\$readmemh}), but the integrated cocotb test builds its own
        load/keep stream rather than consuming that hex. Unifying them is the
        next step.
  \item \textbf{Scale of the head-to-head.} Two small layers (tiny\_cnn,
        mnist\_cnn). They were sufficient to demonstrate divergence but are not
        a full-network evaluation.
  \item \textbf{Energy is derived}, from a per-byte DRAM constant, not from
        switching-activity analysis.
  \item \textbf{FPGA-relative area.} LUT/FF counts are ZCU104 out-of-context
        synthesis, not ASIC $\mu m^2$; they are valid for relative comparison
        between the two schemes, not as absolute silicon area.
\end{enumerate}

% =====================================================================
\section{Future work}

\begin{enumerate}[leftmargin=1.6em,itemsep=4pt]
  \item \textbf{Compress the metadata stream.} The op mix is highly regular, so
        run-length encoding of repeated delta patterns should cut the 128-bit
        per-op cost substantially. Alternatives: stream the metadata in batches
        from DRAM, or use a coarser phase granularity.
  \item \textbf{Close the compute-coupling loop} so a bank-conflict stall
        actually back-pressures the systolic array, making the conflict count a
        \emph{performance} number rather than a memory-side event count.
  \item \textbf{Unify the two op-stream generators} --- have the integrated
        cocotb test consume \file{delta\_ops.hex} directly, closing the last gap
        between the standalone and integrated paths.
  \item \textbf{Conflict-aware allocation.} The bank sweep proves conflicts are
        a deterministic function of the address pattern, so the allocator can be
        extended to choose base addresses that guarantee concurrently-read
        operands land in different banks --- turning $\approx$0 conflicts from
        a happy consequence into a guarantee.
  \item \textbf{Extend the delta model} to simultaneous row-and-column shifts,
        which would raise reported reuse above the current lower bound.
  \item \textbf{Add eviction/replacement to PAGED} so the dynamic baseline is
        realistic at network scale.
  \item \textbf{Full-network evaluation} (ResNet-50, MobileNet) instead of two
        small layers, and larger layers where the STAMP-vs-PAGED byte comparison
        is predicted to invert --- a directly falsifiable prediction.
  \item \textbf{Multi-DNN scratchpad partitioning.} Because placement is static,
        partitioning between concurrent DNNs is a compile-time decision with no
        inter-DNN interference in a translation structure --- a structural
        advantage over a shared page table, worth demonstrating.
  \item \textbf{ASIC synthesis} for absolute area/energy, replacing the derived
        per-byte DRAM energy constant with switching-activity analysis.
\end{enumerate}

% =====================================================================
\section{How to evaluate the results, and why they are correct}

\subsection{The correctness argument in one line}

\begin{keybox}{Three-way independent agreement}
The scheme is implemented \textbf{three times, independently} --- the Python
compiler's own bookkeeping, a separately written cycle-accurate FSM model, and
the SystemVerilog hardware --- and all three produce \textbf{identical} values
for every tracked quantity across 572 operations. A disagreement in any cell
would mean the compiler is describing hardware that does not exist.
\end{keybox}

\begin{center}
\small
\begin{tabular}{@{}lrrrl@{}}
\toprule
\textbf{Quantity} & \textbf{Compiler} & \textbf{FSM model} & \textbf{RTL (xsim)} & \\
\midrule
\op{load} operations   & 134 & 134 & 134 & identical \\
\op{move} operations   & 144 & 144 & 144 & identical \\
\op{keep} operations   & 72  & 72  & 72  & identical \\
\op{alloc} operations  & 127 & 127 & 127 & identical \\
\op{zero} operations   & 95  & 95  & 95  & identical \\
Bytes from DRAM        & 157{,}888 & 157{,}888 & 157{,}888 & identical \\
Bytes moved on-chip    & 184{,}320 & 184{,}320 & 184{,}320 & identical \\
Bytes zero-filled      & 77{,}120  & 77{,}120  & 77{,}120  & identical \\
\bottomrule
\end{tabular}
\end{center}

\noindent The three implementations are \file{stamp\_compiler.py},
\file{pysim/stamp\_ref.py} and \file{stamp\_based\_memory\_controller.sv}. The
RTL reads the compiler's \file{delta\_ops.hex} directly, replays all 127 phase
transitions and 572 operations, and reports \file{stats\_bad\_ops = 0} ---
proving every emitted opcode is one the hardware implements.

\subsection{The seven independent validity checks}

\begin{enumerate}[leftmargin=1.6em,itemsep=4pt]
  \item \textbf{Functional correctness against TensorFlow.} Every STAMP and
        PAGED run is checked against a TF 2.21 golden reference with a 5\,\%
        of full-scale tolerance. Achieved: \textbf{0.013\,\%} (tiny) and
        \textbf{0.009\,\%} (mnist) --- roughly $400\times$ inside tolerance. A
        memory scheme that delivered wrong data would fail here first.
  \item \textbf{Hardware counter vs host bookkeeping.}
        \file{stats\_bytes\_loaded} equals \emph{exactly} $4 \times$ the
        host-computed new-word count (410 words $\to$ 1{,}640\,B;
        862 $\to$ 3{,}448\,B) on both layers. Two independent accountings of
        the same quantity agree to the byte.
  \item \textbf{AXI beat conservation.} Total AXI beats $=$ the layout
        prefetcher's background traffic (identical in both schemes: 1{,}836 /
        10{,}608) \textbf{plus exactly} the delta-load beats (410 / 862) in the
        STAMP runs. The paged backend issues zero AXI by design. Nothing is
        unaccounted for.
  \item \textbf{Page-table conservation.} PAGED hits $+$ misses $=$ exactly the
        replayed footprint touches (1{,}764 / 10{,}608).
  \item \textbf{Bank sweep matches interleaving theory cell by cell.} Every one
        of the 20 measured cells matches what $\text{bank}(a) = a \bmod B$ with
        lowest-port-wins predicts. E.g.\ stride-4 at $B{=}4$: all four ports map
        to bank 0, so 3 of 4 stall every cycle $\to$ $3 \times 200 = 600$
        stalled port-cycles. Measured: 600. $B{=}1$ is 0 by construction and
        measures 0.
  \item \textbf{Directed and software test coverage.} 13 directed RTL tests
        (reset, each opcode in isolation, mixed phases, back-to-back phases,
        reset mid-operation, zero-op phases, busy handshaking, unknown-opcode
        detection, phase base addressing, full stream replay) with 48 passing
        assertions and zero failures; plus 27 software tests for the controller
        and 17 for banking.
  \item \textbf{Regression safety.} The full 19/19 golden regression was re-run
        after \emph{each} change (plumbing, testbench work, PT-write fix): all
        PASS every time, with the other members' results byte-identical (tiny OS
        5{,}808 cycles / 0.013\,\%; WS $8\times8$ 153{,}738; mnist OS 36{,}712 /
        0.009\,\%; all five schedulers \file{all\_pass}). The PT-write broadcast
        cannot affect legacy runs because \file{pt\_write\_en} is never asserted
        there --- confirmed by measurement, not just argument.
\end{enumerate}

\subsection{Internal consistency of the sweep}

$6 \text{ layers} \times 4 \text{ capacities} \times 2 \text{ arrays} = 48$
grid points; \file{sweep\_stamp.csv} holds 34 rows and
\file{capacity\_limits.csv} holds 14 --- they sum exactly to 48, so no point was
silently dropped. Each failure records \emph{why} (``Phase 0 doesn't fit in
on-chip memory'').

\subsection{What would falsify the claims}

\begin{itemize}[leftmargin=1.5em,itemsep=3pt]
  \item If the RTL counters diverged from the compiler's, the scheme would not
        be implementable as described.
  \item If the bank sweep showed conflicts \emph{independent} of the address
        pattern, static placement could not help and the central argument
        collapses.
  \item If the off-chip traffic reduction vanished at larger tile sizes, the
        reuse model would be wrong.
\end{itemize}
All three are directly testable with the harness as built --- which is itself
part of the evidence.

\subsection{Credibility markers to point at}

\begin{itemize}[leftmargin=1.5em,itemsep=3pt]
  \item \textbf{I report the results that go against me} --- PAGED's lower AXI
        beat count, STAMP's higher cycle count, STAMP's $4\times$ register cost
        and 32\,\% lower Fmax.
  \item \textbf{I fixed the competitor's bug before comparing.} The paged
        page-table write gate made ports 1--3 permanently miss; I found it,
        fixed it, and re-ran the entire comparison against the corrected, i.e.\
        stronger, baseline.
  \item \textbf{I removed my own hardcoded numbers.} The previous four-scheme
        comparison was deleted, not defended, and the superseded outputs are
        kept clearly labelled so the change is auditable.
  \item \textbf{Every figure is traceable} to a JSON, CSV or synthesis report
        produced by a tool run, and \file{run\_full\_eval.py} contains no
        performance constants at all.
  \item \textbf{Conservative by construction} --- the delta model under-claims
        reuse on diagonal shifts, and \op{zero} was added specifically to stop
        padding inflating the reported saving.
\end{itemize}

% =====================================================================
\section{Rapid-fire answers to likely examiner questions}

\newcommand{\qa}[2]{\vspace{5pt}\noindent\textbf{\textcolor{stampblue}{Q.}
  #1}\\[2pt]\textbf{A.} #2\vspace{2pt}}

\qa{What exactly is novel? Delta compression is not new.}{
Correct. The novelty is threefold: applying the delta to \emph{scratchpad
layout} in a systolic-array accelerator rather than to data values; eliminating
the runtime tag/translation structure entirely by exploiting a statically known
DNN schedule; and --- what no prior simulator has --- measuring the result
cycle-accurately in RTL against a page-table scheme on identical hardware.}

\qa{Why not just use a bigger scratchpad?}{
Orthogonal and cheaper. A larger scratchpad costs area and leakage; delta
fetching costs a compile-time-sized metadata RAM. And no scratchpad is large
enough for real layers --- the reuse problem persists at every size. The scheme
raises the effective bandwidth of whatever scratchpad you have.}

\qa{Your paged scheme moves fewer AXI beats. Doesn't that defeat your claim?}{
On these layers, yes, and I report it. Page faulting pulls whole 4\,KB pages,
which for a small working set is coarse but efficient. My claim is fewer bank
conflicts ($42\times$), no translation hardware, and deterministic placement.
For larger layers where a fault pulls far more than the delta, the byte
comparison inverts --- a falsifiable prediction.}

\qa{How do you know the compiler and the hardware agree?}{
Three independent implementations produce identical values for all eight tracked
quantities across 572 operations, and the RTL reads the compiler's hex file
directly with \file{stats\_bad\_ops = 0}.}

\qa{What does \op{zero} do and why does it exist?}{
Padded convolutions have receptive fields extending past the tensor border.
Those bytes are implicit zeros: present on-chip, never fetched. Before the
opcode existed the compiler charged them as DRAM loads, overstating traffic by
77{,}120\,B. Adding it reduced the reported load volume from 235{,}008 to
157{,}888\,B --- it made the number \emph{worse for the story and truer}.}

\qa{What is the hardware cost of your scheme?}{
A nine-state FSM plus the metadata RAM. The tag array, comparators and
replacement logic disappear. Measured at 256 entries: 6{,}434 LUTs / 22{,}198
FFs at 405\,MHz, versus 6{,}110 / 5{,}440 at 594\,MHz for the paged manager ---
so similar LUTs, $4\times$ the registers, 32\,\% lower Fmax. The metadata RAM
(9{,}152\,B for the 128-phase layer) is the real cost and the main scaling limit.}

\qa{Why do the two schemes differ so much in bank conflicts?}{
The stamp compiler picks physical addresses at compile time and bump-allocates
densely, so concurrent reads are stride-1 and land in different banks. The page
table assigns physical pages at runtime and preserves page-relative offsets, so
row-strided reads collide. The bank sweep proves conflicts are a deterministic
function of the address pattern --- exactly what static placement controls.}

\qa{Isn't your bank-conflict estimate in Experiment 3 just a formula?}{
Yes, and it is labelled as such. It is a birthday-problem bound used to cover 14
workloads where full RTL simulation is impractical. Experiment 4b is the
authority --- cycle-exact counters. The measured values sit far below the bound
because real access patterns are consecutive, not random, and that gap is itself
a finding.}

\qa{How would this behave on a multi-DNN workload?}{
Each DNN gets its own stamp schedule and the scratchpad is partitioned at
compile time. Because placement is static there is no inter-DNN interference in
a translation structure --- an advantage over a shared page table needing
multiple coexisting address spaces. Demonstrating it is future work.}

\end{document}
